\documentclass[11pt]{article}

\usepackage[utf8]{inputenc}
\usepackage[T1]{fontenc}
\usepackage{lmodern}
\usepackage{microtype}
\usepackage{amsmath,amssymb}
\usepackage{booktabs}
\usepackage{array}
\usepackage{xcolor}
\usepackage[margin=1in]{geometry}
\usepackage{enumitem}
\usepackage{titlesec}
\usepackage{titling}
\usepackage[hidelinks]{hyperref}

% ---- Palette (matched to the 7 July proposal) ------------------------------
\definecolor{ink}{RGB}{20,25,40}
\definecolor{accent}{RGB}{27,54,93}      % sober navy
\definecolor{rule}{RGB}{200,205,215}
\definecolor{softgrey}{RGB}{95,100,110}
\definecolor{gooddk}{RGB}{27,94,60}      % done
\definecolor{pend}{RGB}{140,90,20}       % pending

\hypersetup{colorlinks=true, linkcolor=accent, urlcolor=accent, citecolor=accent}

\titleformat{\section}
  {\color{accent}\normalfont\large\bfseries}{\thesection.}{0.6em}{}
\titleformat{\subsection}
  {\color{ink}\normalfont\normalsize\bfseries}{\thesubsection}{0.6em}{}
\titlespacing*{\section}{0pt}{1.4em}{0.5em}
\titlespacing*{\subsection}{0pt}{1.0em}{0.35em}

\setlist[itemize]{leftmargin=1.4em, itemsep=0.2em, topsep=0.25em}
\setlist[enumerate]{leftmargin=1.6em, itemsep=0.25em, topsep=0.25em}

\newcommand{\done}{\textcolor{gooddk}{\textbf{Done}}}
\newcommand{\ready}{\textcolor{gooddk}{Ready}}
\newcommand{\pending}{\textcolor{pend}{Pending GPU}}

\setlength{\parskip}{0.45em}
\setlength{\parindent}{0pt}

% ---------------------------------------------------------------------------
\begin{document}

% ===== Title block =========================================================
\begin{center}
{\color{accent}\LARGE\bfseries From Diagnosis to Solution}\\[0.35em]
{\large SLEEP: Implementation Milestone and Validation Plan}\\[0.6em]
{\color{softgrey}\normalsize Interim progress report, ahead of empirical validation}
\end{center}

\vspace{0.4em}
{\color{rule}\hrule height 0.8pt}
\vspace{0.6em}

\begin{center}
\begin{tabular}{@{}r l@{\hspace{3em}}r l@{}}
\textbf{Prepared for:} & Prof.\ Debashis Guha & \textbf{Date:} & 27 July 2026 \\
 & Director, MAIB; Chair, CRTIB & \textbf{Stage:} & Interim (2 of 3) \\
 & SP Jain School of Global Management & \textbf{Follows:} & Proposal, 7 July 2026 \\[0.5em]
\textbf{Prepared by:} & Aditya Tripathi & & \\
\end{tabular}
\end{center}

\vspace{0.5em}
{\color{rule}\hrule height 0.8pt}
\vspace{0.8em}

% ===== Preamble ============================================================
Professor Guha,

This interim report covers the three weeks since the proposal of 7 July. The
engineering phase is complete: every code change the revision plan called for
--- the three-phase P1/P2/P3 roadmap and the retrieval-aware warm-up extension
--- is implemented and verified, with the full test suite green (\textbf{362
tests passing, 0 failing}, up from 297 before this phase). What remains is the
GPU validation on RunPod, whose numbers will populate the final report. The
report is deliberately honest about the boundary between the two: the code that
\emph{could} close the recognition--recall gap now exists and is correct; whether
it \emph{does} is the empirical question the runs will answer. Below: what was
built, how it was verified, and the runs that will produce the results.

% ===== 1. Progress at a glance =============================================
\section{Progress against the plan}

Every item from the proposal's roadmap is implemented and unit-tested on CPU
(via tiny in-memory models), or, for the model/precision variants, staged as
run-ready configuration. Nothing has yet been executed at 7B scale.

\newcolumntype{L}{>{\raggedright\arraybackslash}p{8.4cm}}
\begin{center}
\small
\begin{tabular}{@{}L l@{}}
\toprule
\multicolumn{2}{@{}l}{\textbf{P1 --- Critical}} \\
\midrule
Multi-seed harness ($\pm$SD aggregation) + seeding wired through the pipeline & \done \\
Extended multi-cycle (10--20 cycles via existing cycle/batch controls) & \ready \\
Multi-model \& multi-precision configs (Llama~3, Mistral, Qwen~1.5B, float32) & \ready \\
\midrule
\multicolumn{2}{@{}l}{\textbf{P2 --- High}} \\
\midrule
Standard-benchmark loaders (LAMA, CounterFact, real post-cutoff facts) & \done \\
Missing baselines: RAG, EWC-only, in-context injection & \done \\
Two-stage validation (surprise pre-filter + direct mini-recall gate) & \done \\
Biological-framing revision (constrain-or-trim) & Paper edit \\
\midrule
\multicolumn{2}{@{}l}{\textbf{P3 --- Polish \& Solution}} \\
\midrule
\textbf{Retrieval-aware warm-up extension} (gate + trainer + injector seam) & \done \\
Calibration plot (confidence vs.\ accuracy, tagged vs.\ untagged) & \done \\
Failure-mode table (C1--C5) & Paper edit \\
\bottomrule
\end{tabular}
\end{center}

``Done'' means implemented and covered by passing tests; ``Ready'' means the
code path is exercised by tests and the run is a matter of configuration on the
pod; ``Paper edit'' items are writing tasks folded into the final report. The
two model/precision items are ``Ready'' rather than ``Done'' because the
component paths are tested but the specific external models have not been loaded.

% ===== 2. The extension ====================================================
\section{The extension is built --- and verified where it can be}

The centrepiece of the proposal was the retrieval-aware warm-up: a small
trainable gate that teaches a frozen model to \emph{use} key--value (KV) memory
during generation, addressing the recognition--recall gap the paper diagnoses.
It is now implemented as three pieces:

\begin{itemize}
  \item A \textbf{memory gate} --- a per-layer, per-head multiplicative scale on
    the memory value contribution (roughly 36 parameters for the top third of
    Qwen2.5-7B). Because attention is linear in the values, this exactly
    modulates how strongly retrieved memory feeds generation. It is initialised
    to identity, so an untrained gate changes nothing.
  \item A \textbf{warm-up trainer} that, on a general corpus (never the
    evaluation facts), stores a passage's prefix in memory and trains the gate
    to predict the continuation from it --- the Memorizing-Transformers-style
    objective, teaching the \emph{skill} of using memory, not any fact.
  \item An \textbf{injector seam} so the gate lives on the attention path rather
    than in the consolidation adapter, and therefore persists across sleep
    cycles instead of being overwritten by them.
\end{itemize}

Crucially, the change touches the KV-injection critical path, so it was verified
against ground truth before anything was built on top of it. Per our standing
practice, the \emph{first} test written was a hook-based equality check: with an
identity gate installed, the injected-attention output must be bit-identical to
the ungated path. It is (to $10^{-6}$). The 54 pre-existing KV-injection and
memory tests continue to pass unchanged, and the gate is confirmed
differentiable end-to-end so the warm-up can actually train it. In short, the
mechanism is in place and provably does not perturb the verified recognition
pathway until it is deliberately trained.

While there, we generalised the injector to resolve a model's attention module
dynamically, so KV injection nominally runs on Llama and Mistral as well as
Qwen. This is what makes the cross-model recognition test (P1) feasible; we flag
in Section~\ref{sec:risks} that only the Qwen path is bit-verified so far.

% ===== 3. Rigor ============================================================
\section{Verification and rigor}

The engineering is not asserted, it is tested. Concretely:

\begin{itemize}
  \item \textbf{362 tests pass, 0 fail} --- about 65 new tests over the prior
    297, covering the gate, the warm-up loop end-to-end on a tiny in-memory
    model, the seeding utilities, the benchmark loaders, the two new baselines,
    the calibration plot, and the two-stage-validation gate.
  \item \textbf{Ground-truth-first on the critical path.} The identity-gate
    equality check and a gate-gradient check were written before the trainer,
    so a silent change to injection could not masquerade as a warm-up effect.
  \item \textbf{The self-grading gap is now instrumented.} Two-stage validation
    records how many candidates pass the surprise proxy versus the direct
    free-form recall gate, so the C4 discrepancy the paper names is measured as
    a number rather than argued.
  \item \textbf{Reproducibility built in.} A single call seeds every RNG the
    pipeline touches, and the multi-seed runner executes an experiment across
    seeds in independent processes and reports mean~$\pm$~standard deviation ---
    directly answering the review's ``no single-seed results'' requirement.
\end{itemize}

Everything above runs on CPU, so correctness was established without spending
GPU budget. The 7B runs are for measurement, not debugging.

% ===== 4. What the runs will decide ========================================
\section{What the validation runs will decide}
\label{sec:runs}

The implementation cannot, by itself, tell us whether the extension works ---
only whether it is correct. The scientific claims remain open until the runs
land, and the report is explicit about that. The pre-registered success
criterion is unchanged from the proposal: if warmed-up SLEEP reaches
$\mathrm{DRA}\approx0.10$--$0.15$ at $\mathrm{BCP}<1.5$ on the 200-fact
benchmark, the recognition--recall gap is closable within the framework. The
runs, in priority order and with the number each produces:

\begin{enumerate}
  \item \textbf{Multi-seed recognition \& Pareto} (seeds 0--4): does the
    $+0.16$ recognition signal and the single-cycle frontier survive seed
    variance, as mean~$\pm$~SD?
  \item \textbf{Extended multi-cycle} (10--20 cycles, seeds 0--4): does base-
    capability preservation plateau or keep climbing beyond three cycles?
  \item \textbf{Warm-up extension} --- the headline: the SLEEP-with-warm-up row
    of the new results table (DRA, BCP) against the without-warm-up baseline.
  \item \textbf{float32 and a second model family}: is the C1 precision-floor
    finding a bfloat16 artefact, and does the gap generalise beyond Qwen?
  \item \textbf{Standard benchmark \& baselines}: the three baselines and an
    established dataset, for comparability with prior work.
\end{enumerate}

Estimated compute is roughly USD~\$8--15 depending on the seed and cycle counts,
within the pod budget. Each run writes a machine-readable result file, so the
final report is a matter of tabulating outputs, not re-deriving them.

% ===== 5. Risks ============================================================
\section{Open risks}
\label{sec:risks}

Three, stated plainly:

\begin{itemize}
  \item \textbf{The extension may only narrow, not close, the gap.} A partial
    rise in DRA is still a publishable diagnostic-to-treatment result, but the
    magnitude is unknown until the run. We are not assuming the modest success
    criterion will be met.
  \item \textbf{Cross-model KV injection is CPU-verified for Qwen only.} The
    Llama/Mistral attention path is generalised but not yet bit-checked on
    hardware; the first cross-model run doubles as that check, and the
    lower-risk float32 and multi-cycle comparisons do not depend on it.
  \item \textbf{Compute state.} The pod's connection details and remaining
    credit change between sessions; we will confirm both before starting, and
    lead with the cheapest high-signal run (the warm-up on a 50-fact subset,
    about \$1) as a smoke test before committing to the full seed sweeps.
\end{itemize}

% ===== 6. Timeline =========================================================
\section{Revised timeline}

\begin{center}
\small
\begin{tabular}{@{}l l l@{}}
\toprule
Stage & Deliverable & Status \\
\midrule
1. Proposal        & Roadmap + extension, sign-off request & \done\ (7 Jul) \\
2. Implementation  & All code + tests; this report          & \done\ (27 Jul) \\
3. Validation      & GPU runs on RunPod (\S\ref{sec:runs})   & \pending\ ($\sim$1 week) \\
4. Final report    & Results, tables, paper revision         & After runs \\
\bottomrule
\end{tabular}
\end{center}

The project is on the plan submitted three weeks ago. The build phase came in
where scoped; the only work between here and the final report is executing the
runs and writing up what they show.

% ===== 7. Closing ==========================================================
\section{Requested input}

No decision is required to proceed --- the runs are queued and can begin as soon
as the pod is available. Two light confirmations would help us aim the final
report:

\begin{enumerate}
  \item That the validation priority order in Section~\ref{sec:runs} matches how
    you would weight the evidence, or your steer if not.
  \item A target venue and deadline, so the final report is written to the right
    format and length from the start.
\end{enumerate}

We will report the smoke-test result as soon as it is in, and the full tables in
the final report to follow.

\vspace{0.8em}
{\color{rule}\hrule height 0.6pt}
\vspace{0.4em}
{\color{softgrey}\small Respectfully submitted --- Aditya Tripathi. Implementation, tests, and experiment scripts are in the project repository; the 7 July proposal accompanies this report.}

\end{document}
